% ==============================================================================
% Plantilla Institucional de Trabajo Terminal — UPIIZ IPN
% Archivo: capitulos/protocolo/04-marco-teorico.tex
% ==============================================================================
% ESTRUCTURA NORMATIVA DE PROTOCOLO: CAPÍTULO 4 - MARCO TEÓRICO
% ==============================================================================
% LÍMITE INSTITUCIONAL: Máximo 3 páginas.
%
% Requisito normativo (Anexo 1):
% "Presentar los antecedentes teóricos que constituyen el sustento del trabajo
% a desarrollar, utilizando paráfrasis y referencias correctas."
%
% Criterio de aplicación:
% - El contenido específico dependerá exclusivamente de la naturaleza del proyecto
%   (mecánica, electrónica, control, software, visión artificial, automatización, etc.).
% - No asumir como obligatorios temas particulares como modelado dinámico o teoría
%   de control si no corresponden a la solución propuesta.
% ==============================================================================

\chapter{Marco Teórico}
\label{cap:proto-marco-teorico}

% [INSTRUCCIÓN PARA EL ALUMNO]:
% Desarrolle en un máximo de 3 páginas los conceptos, principios, modelos analíticos
% o fundamentos conceptuales que sustentan técnicamente el trabajo a desarrollar.
%
% Recuerde redactar utilizando rigurosamente la técnica de paráfrasis e incorporar
% las referencias bibliográficas correspondientes en formato IEEE (ej. \cite{...}).

\section{Fundamentos teóricos del proyecto}

[Exponga los fundamentos teóricos indispensables que respaldan la propuesta, seleccionando aquellos temas que guían directamente las decisiones de diseño o desarrollo del sistema].

\section{Modelos y principios aplicables a la solución}

[Presente las relaciones analíticas, modelos conceptuales, arquitecturas lógicas o principios de operación necesarios para el desarrollo del proyecto, adaptados a la naturaleza del sistema propuesto].
